% ============================================================
% 培优专题13 — 构造法与存在性问题
% 编译命令：xelatex -interaction=nonstopmode 培优专题13_构造法与存在性问题.tex
% ============================================================
\documentclass[12pt, a4paper]{ctexart}
\usepackage[margin=2.2cm]{geometry}
\usepackage{amsmath, amssymb, amsfonts}
\usepackage{enumitem}
\usepackage{tcolorbox}
\usepackage{booktabs}
\usepackage{array}
\usepackage{fancyhdr}
\usepackage{xcolor}
\usepackage{tikz}
\usepackage{hyperref}
\usepackage{graphicx}
\usepackage{multirow}
\usepackage{pgfplots}
\usepackage{amsthm}
\pgfplotsset{compat=1.18}

% 颜色定义
\definecolor{mainblue}{RGB}{41, 98, 168}
\definecolor{lightblue}{RGB}{230, 240, 255}
\definecolor{tipgreen}{RGB}{34, 139, 87}
\definecolor{lightgreen}{RGB}{235, 250, 240}
\definecolor{warnred}{RGB}{200, 50, 50}
\definecolor{lightred}{RGB}{255, 240, 240}
\definecolor{examorange}{RGB}{210, 120, 20}
\definecolor{lightorange}{RGB}{255, 248, 235}
\definecolor{apppurple}{RGB}{120, 60, 170}
\definecolor{lightpurple}{RGB}{245, 238, 255}

% 页眉页脚
\pagestyle{fancy}
\fancyhf{}
\fancyhead[L]{\small\textcolor{mainblue}{\textbf{高中数学培优 $\cdot$ 构造法与存在性问题}}}
\fancyhead[R]{\small\textcolor{mainblue}{培优专题13}}
\fancyfoot[C]{\small —\ \thepage\ —}
\renewcommand{\headrulewidth}{0.6pt}
\renewcommand{\headrule}{\hbox to\headwidth{\color{mainblue}\leaders\hrule height \headrulewidth\hfill}}
\renewcommand{\footrulewidth}{0pt}

% tcolorbox 样式
\newtcolorbox{knowledgebox}[1][]{%
  colback=lightblue, colframe=mainblue,
  fonttitle=\bfseries\large, title=#1,
  arc=3mm, boxrule=0.8pt,
  left=8pt, right=8pt, top=6pt, bottom=6pt
}

\newtcolorbox{tipbox}[1][]{%
  colback=lightgreen, colframe=tipgreen,
  fonttitle=\bfseries, title=#1,
  arc=2mm, boxrule=0.6pt,
  left=6pt, right=6pt, top=4pt, bottom=4pt
}

\newtcolorbox{warnbox}[1][]{%
  colback=lightred, colframe=warnred,
  fonttitle=\bfseries, title=#1,
  arc=2mm, boxrule=0.6pt,
  left=6pt, right=6pt, top=4pt, bottom=4pt
}

\newtcolorbox{examplebox}[1][]{%
  colback=lightorange, colframe=examorange,
  fonttitle=\bfseries, title=#1,
  arc=3mm, boxrule=0.8pt,
  left=8pt, right=8pt, top=6pt, bottom=6pt
}

\newtcolorbox{appbox}[1][]{%
  colback=lightpurple, colframe=apppurple,
  fonttitle=\bfseries, title=#1,
  arc=3mm, boxrule=0.8pt,
  left=8pt, right=8pt, top=6pt, bottom=6pt
}

\setCJKmainfont{Hiragino Sans GB}
\setCJKsansfont{Heiti SC}

\begin{document}

% ============ 标题页 ============
\begin{center}
{\Huge\bfseries\textcolor{mainblue}{培优专题13\quad 构造法与存在性问题}}\\[12pt]
{\Large\textcolor{gray}{培\quad 优\quad 讲\quad 义}}\\[8pt]
\rule{0.6\textwidth}{1pt}\\[8pt]
{\normalsize 适用：高二高三拔高\quad |\quad 课时：3课时}
\end{center}
\vspace{10pt}

% ============ 第一部分：知识点讲解 ============
\section*{\textcolor{mainblue}{一、知识点讲解}}

\subsection*{1.1\quad 构造函数证明不等式}

\begin{knowledgebox}[核心思想]
构造函数法是将不等式问题转化为函数问题的基本策略。其核心在于：\textbf{将待证不等式的某一边视为零}，构造一个新函数，利用函数的单调性、极值等性质来证明不等式。

\medskip
\textbf{基本步骤：}
\begin{enumerate}[label=(\arabic*)]
  \item 将不等式 $f(x) > g(x)$ 改写为 $F(x) = f(x) - g(x) > 0$；
  \item 分析 $F(x)$ 的定义域、单调性、极值；
  \item 利用 $F(x)_{\min} \geq 0$ 或 $F(x)_{\max} \leq 0$ 完成证明。
\end{enumerate}
\end{knowledgebox}

\begin{tipbox}[常用构造技巧]
\textbf{技巧一：移项构造}\\
若证 $f(x) \geq g(x)$，令 $F(x) = f(x) - g(x)$，证 $F(x) \geq 0$。

\medskip
\textbf{技巧二：作商构造}\\
若 $g(x) > 0$，令 $H(x) = \dfrac{f(x)}{g(x)}$，证 $H(x) \geq 1$。

\medskip
\textbf{技巧三：取对数构造}\\
若 $f(x) > 0,\, g(x) > 0$，证 $f(x) \geq g(x)$ 等价于证 $\ln f(x) \geq \ln g(x)$。

\medskip
\textbf{技巧四：分离参数构造}\\
将含参不等式分离参数，如 $a \geq h(x)$，转化为求 $h(x)$ 的最大值。
\end{tipbox}

\subsection*{1.2\quad 构造辅助数列}

\begin{knowledgebox}[递推数列的构造]
当数列 $\{a_n\}$ 满足复杂的递推关系时，可通过构造辅助数列将递推关系化简。

\medskip
\textbf{常见构造类型：}
\begin{enumerate}[label=(\arabic*)]
  \item \textbf{等差型构造}：若 $a_{n+1} = a_n + f(n)$，累加可得通项。
  \item \textbf{等比型构造}：若 $a_{n+1} = pa_n + q$，令 $b_n = a_n + \dfrac{q}{p-1}$，则 $\{b_n\}$ 为等比数列。
  \item \textbf{取倒数构造}：若 $a_{n+1} = \dfrac{a_n}{1 + ca_n}$，令 $b_n = \dfrac{1}{a_n}$ 化简。
  \item \textbf{取对数构造}：若 $a_{n+1} = a_n^k$（$a_n > 0$），令 $b_n = \ln a_n$ 化为线性递推。
  \item \textbf{特征根构造}：对于 $a_{n+2} = pa_{n+1} + qa_n$，利用特征方程 $x^2 = px + q$ 的根构造。
\end{enumerate}
\end{knowledgebox}

\subsection*{1.3\quad 反例与否定策略}

\begin{knowledgebox}[反例法]
\textbf{反例法}是通过构造一个满足条件但不满足结论的例子来推翻一个命题。

\medskip
\textbf{适用场景：}
\begin{itemize}[label=$\bullet$]
  \item 判断命题的真假（特别是"对任意""所有"类型的全称命题）；
  \item 否定一个猜想；
  \item 在选择题中快速排除错误选项。
\end{itemize}

\medskip
\textbf{构造反例的原则：}
\begin{enumerate}[label=(\arabic*)]
  \item 从简单情形入手（特殊值、小数值）；
  \item 利用极端情形（边界值、趋近于 $0$ 或 $+\infty$）；
  \item 利用对称性和周期性简化讨论；
  \item 逐条检验条件，确保反例满足全部前提。
\end{enumerate}
\end{knowledgebox}

\subsection*{1.4\quad 存在性问题的策略}

\begin{knowledgebox}[存在性问题]
\textbf{存在性问题}是指证明存在某个对象（数、函数、点等）满足特定条件。

\medskip
\textbf{常用策略：}
\begin{enumerate}[label=(\arabic*)]
  \item \textbf{直接构造法}：直接构造出满足条件的对象，这是最有力的证明方式。
  \item \textbf{反证法}：假设不存在，推导出矛盾。
  \item \textbf{中间值定理法}：若 $f$ 在 $[a,b]$ 上连续且 $f(a) \cdot f(b) < 0$，则 $\exists\, c \in (a,b)$ 使 $f(c) = 0$。
  \item \textbf{抽屉原理}：$n+1$ 个物体放入 $n$ 个抽屉，必有一个抽屉不少于 $2$ 个。
  \item \textbf{极端原理}：在有限集合中，最大值和最小值一定存在。
  \item \textbf{数学归纳法}：适用于与自然数相关的存在性问题。
\end{enumerate}
\end{knowledgebox}

\begin{warnbox}[易错提醒]
\begin{enumerate}[label=(\arabic*)]
  \item "存在"只需找到一个即可，不要求找出所有满足条件的对象。
  \item 反证法假设时要注意否定的正确性："存在 $x$ 使得 $P(x)$"的否定是"对所有 $x$，$\lnot P(x)$"。
  \item 中间值定理要求函数\textbf{连续}，不能遗漏此前提。
  \item 构造反例时一定要逐一验证所有前提条件。
\end{enumerate}
\end{warnbox}

% ============ 第二部分：例题精讲 ============
\section*{\textcolor{examorange}{二、例题精讲}}

\begin{examplebox}[【例题1】构造函数证明经典不等式]
\textbf{题目：}已知 $x > 0$，证明：$\mathrm{e}^x > 1 + x + \dfrac{x^2}{2}$。

\medskip
\textbf{证明：}

设 $f(x) = \mathrm{e}^x - 1 - x - \dfrac{x^2}{2}$，则

$$f'(x) = \mathrm{e}^x - 1 - x$$

令 $g(x) = f'(x) = \mathrm{e}^x - 1 - x$，则

$$g'(x) = \mathrm{e}^x - 1$$

当 $x > 0$ 时，$g'(x) = \mathrm{e}^x - 1 > 0$，所以 $g(x)$ 在 $(0, +\infty)$ 上单调递增。

因此当 $x > 0$ 时，$g(x) > g(0) = 0$，即 $f'(x) > 0$。

从而 $f(x)$ 在 $(0, +\infty)$ 上单调递增。

当 $x > 0$ 时，$f(x) > f(0) = 1 - 1 - 0 - 0 = 0$。

故 $\mathrm{e}^x > 1 + x + \dfrac{x^2}{2}$ 对所有 $x > 0$ 成立。\hfill$\square$
\end{examplebox}

\begin{examplebox}[【例题2】构造函数证明对数不等式]
\textbf{题目：}已知 $a > b > 0$，证明：$\ln a - \ln b > \dfrac{2(a-b)}{a+b}$。

\medskip
\textbf{证明：}

令 $t = \dfrac{a}{b} > 1$，则 $a = tb$，不等式化为：

$$\ln t > \frac{2(tb - b)}{tb + b} = \frac{2(t-1)}{t+1}$$

设 $F(t) = \ln t - \dfrac{2(t-1)}{t+1}$（$t > 1$），则

$$F'(t) = \frac{1}{t} - \frac{2(t+1) - 2(t-1)}{(t+1)^2} = \frac{1}{t} - \frac{4}{(t+1)^2} = \frac{(t+1)^2 - 4t}{t(t+1)^2} = \frac{(t-1)^2}{t(t+1)^2}$$

当 $t > 1$ 时，$F'(t) = \dfrac{(t-1)^2}{t(t+1)^2} > 0$。

因此 $F(t)$ 在 $(1, +\infty)$ 上单调递增。

当 $t > 1$ 时，$F(t) > F(1) = \ln 1 - 0 = 0$。

故 $\ln a - \ln b > \dfrac{2(a-b)}{a+b}$ 成立。\hfill$\square$
\end{examplebox}

\begin{tipbox}[方法总结]
\textbf{换元简化}是处理多变量不等式的常用手段。通过 $t = a/b$ 将二元问题转化为一元问题，大大降低了构造函数求导的难度。对于涉及 $\ln a - \ln b$ 和 $a - b$ 的不等式，此法尤为有效。
\end{tipbox}

\begin{examplebox}[【例题3】构造辅助数列求通项]
\textbf{题目：}已知数列 $\{a_n\}$ 满足 $a_1 = 1$，$a_{n+1} = 2a_n + 3$，求通项公式 $a_n$。

\medskip
\textbf{解：}

设 $a_{n+1} + \lambda = 2(a_n + \lambda)$，展开得 $a_{n+1} = 2a_n + \lambda$。

对比 $a_{n+1} = 2a_n + 3$，得 $\lambda = 3$。

令 $b_n = a_n + 3$，则 $b_{n+1} = 2b_n$，$\{b_n\}$ 是公比为 $2$ 的等比数列。

$$b_1 = a_1 + 3 = 4$$

$$b_n = 4 \cdot 2^{n-1} = 2^{n+1}$$

$$\boxed{a_n = 2^{n+1} - 3}$$
\end{examplebox}

\begin{examplebox}[【例题4】取倒数构造化简递推]
\textbf{题目：}已知 $a_1 = 1$，$a_{n+1} = \dfrac{a_n}{1 + 2a_n}$，求通项公式 $a_n$。

\medskip
\textbf{解：}

对递推式两边取倒数：

$$\frac{1}{a_{n+1}} = \frac{1 + 2a_n}{a_n} = \frac{1}{a_n} + 2$$

令 $b_n = \dfrac{1}{a_n}$，则 $b_{n+1} = b_n + 2$，$\{b_n\}$ 是公差为 $2$ 的等差数列。

$$b_1 = \frac{1}{a_1} = 1$$

$$b_n = 1 + 2(n-1) = 2n - 1$$

$$\boxed{a_n = \frac{1}{2n-1}}$$
\end{examplebox}

\begin{examplebox}[【例题5】构造反例否定命题]
\textbf{题目：}判断以下命题的真假，并说明理由。

\textbf{命题P}：若函数 $f(x)$ 在 $(a,b)$ 上可导且 $f'(x) > 0$，则 $f(x)$ 在 $(a,b)$ 上恒为正。

\medskip
\textbf{解：}

\textbf{命题P为假。}

\textbf{反例：}取 $f(x) = x$，定义域为 $(-1, 1)$。

则 $f'(x) = 1 > 0$，满足条件。

但 $f(-0.5) = -0.5 < 0$，不满足"$f(x)$ 恒为正"。

因此命题P为假命题。\hfill$\square$

\medskip
\textbf{反例构造思路：}单调递增只能保证函数值"越来越大"，但无法保证函数值为正。最简单的反例就是过原点的直线。
\end{examplebox}

\begin{examplebox}[【例题6】构造函数证明双变量不等式]
\textbf{题目：}已知 $x, y > 0$ 且 $x + y = 1$，证明：$\left(1 + \dfrac{1}{x}\right)\left(1 + \dfrac{1}{y}\right) \geq 9$。

\medskip
\textbf{证明：}

将 $y = 1 - x$（$0 < x < 1$）代入，设

$$f(x) = \left(1 + \frac{1}{x}\right)\left(1 + \frac{1}{1-x}\right) = \frac{(x+1)(2-x)}{x(1-x)}$$

化简分子：$(x+1)(2-x) = -x^2 + x + 2$。

化简分母：$x(1-x) = x - x^2$。

$$f(x) = \frac{-x^2 + x + 2}{x - x^2} = 1 + \frac{2}{x - x^2} = 1 + \frac{2}{x(1-x)}$$

由 AM-GM 不等式，$x(1-x) \leq \left(\dfrac{x + (1-x)}{2}\right)^2 = \dfrac{1}{4}$，

等号当 $x = \dfrac{1}{2}$ 时成立。

因此 $f(x) = 1 + \dfrac{2}{x(1-x)} \geq 1 + \dfrac{2}{\frac{1}{4}} = 1 + 8 = 9$。

等号当 $x = y = \dfrac{1}{2}$ 时成立。\hfill$\square$
\end{examplebox}

\begin{examplebox}[【例题7】特征根法构造辅助数列]
\textbf{题目：}已知数列 $\{a_n\}$ 满足 $a_1 = 1$，$a_2 = 2$，$a_{n+2} = 3a_{n+1} - 2a_n$，求通项公式 $a_n$。

\medskip
\textbf{解：}

特征方程：$x^2 = 3x - 2$，即 $x^2 - 3x + 2 = 0$。

解得 $x_1 = 1$，$x_2 = 2$。

设 $a_n = A \cdot 1^n + B \cdot 2^n = A + B \cdot 2^n$。

由初始条件：
$$\begin{cases} a_1 = A + 2B = 1 \\ a_2 = A + 4B = 2 \end{cases}$$

两式相减得 $2B = 1$，$B = \dfrac{1}{2}$，$A = 0$。

$$\boxed{a_n = 2^{n-1}}$$
\end{examplebox}

\begin{examplebox}[【例题8】构造函数证明带参数不等式]
\textbf{题目：}已知 $f(x) = x - \ln(1+x)$，若对所有 $x \geq 0$ 都有 $f(x) \leq ax^2$，求实数 $a$ 的取值范围。

\medskip
\textbf{解：}

当 $x = 0$ 时，$f(0) = 0 = a \cdot 0$，对任意 $a$ 成立。

当 $x > 0$ 时，$f(x) \leq ax^2$ 等价于 $a \geq \dfrac{x - \ln(1+x)}{x^2}$。

设 $g(x) = \dfrac{x - \ln(1+x)}{x^2}$（$x > 0$），需求 $g(x)$ 的最大值。

$$g'(x) = \frac{\left(1 - \frac{1}{1+x}\right) x^2 - [x - \ln(1+x)] \cdot 2x}{x^4} = \frac{\frac{x^3}{1+x} - 2x[x - \ln(1+x)]}{x^4}$$

$$= \frac{\frac{x}{1+x} - 2 + \frac{2\ln(1+x)}{x}}{x^2} \cdot x = \frac{x^2 - 2(x+1)[x - \ln(1+x)]}{x^3(1+x)}$$

设 $h(x) = x^2 - 2(x+1)[x - \ln(1+x)] = x^2 - 2x^2 - 2x + 2(x+1)\ln(1+x)$

$= -x^2 - 2x + 2(x+1)\ln(1+x)$。

$h'(x) = -2x - 2 + 2\ln(1+x) + 2 = -2x + 2\ln(1+x)$。

$h''(x) = -2 + \dfrac{2}{1+x} = \dfrac{-2x}{1+x} < 0$（$x > 0$）。

所以 $h'(x)$ 在 $(0, +\infty)$ 上单调递减，$h'(x) < h'(0) = 0$。

因此 $h(x)$ 在 $(0, +\infty)$ 上单调递减，$h(x) < h(0) = 0$。

从而 $g'(x) < 0$（$x > 0$），$g(x)$ 在 $(0, +\infty)$ 上单调递减。

由洛必达法则：

$$\lim_{x \to 0^+} g(x) = \lim_{x \to 0^+} \frac{x - \ln(1+x)}{x^2} = \lim_{x \to 0^+} \frac{1 - \frac{1}{1+x}}{2x} = \lim_{x \to 0^+} \frac{x}{2x(1+x)} = \frac{1}{2}$$

故 $g(x) < \dfrac{1}{2}$（$x > 0$），即 $\dfrac{x - \ln(1+x)}{x^2} < \dfrac{1}{2}$。

因此 $a \geq \dfrac{1}{2}$ 时不，$f(x) \leq ax^2$ 对所有 $x \geq 0$ 成立。

$$\boxed{a \geq \frac{1}{2}}$$
\end{examplebox}

\begin{examplebox}[【例题9】反证法证明存在性]
\textbf{题目：}证明：方程 $x^3 + x + 1 = 0$ 至少有一个实数根。

\medskip
\textbf{证明：}

设 $f(x) = x^3 + x + 1$，则 $f(x)$ 在 $\mathbb{R}$ 上连续。

$$f(-1) = -1 - 1 + 1 = -1 < 0$$

$$f(0) = 0 + 0 + 1 = 1 > 0$$

由中间值定理，$\exists\, c \in (-1, 0)$ 使得 $f(c) = 0$。

故方程 $x^3 + x + 1 = 0$ 在区间 $(-1, 0)$ 内至少有一个实数根。\hfill$\square$

\medskip
\textbf{进一步讨论：}$f'(x) = 3x^2 + 1 > 0$，故 $f(x)$ 在 $\mathbb{R}$ 上严格单调递增，方程恰好有一个实数根。
\end{examplebox}

\begin{examplebox}[【例题10】构造函数证明数列不等式]
\textbf{题目：}证明：$\displaystyle\sum_{k=1}^{n} \frac{1}{k^2} < 2 - \frac{1}{n}$（$n \geq 2$）。

\medskip
\textbf{证明：}

注意到当 $k \geq 2$ 时，$\dfrac{1}{k^2} < \dfrac{1}{k(k-1)} = \dfrac{1}{k-1} - \dfrac{1}{k}$。

因此当 $n \geq 2$ 时：

$$\sum_{k=1}^{n} \frac{1}{k^2} = 1 + \sum_{k=2}^{n} \frac{1}{k^2} < 1 + \sum_{k=2}^{n} \left(\frac{1}{k-1} - \frac{1}{k}\right)$$

裂项求和：

$$\sum_{k=2}^{n} \left(\frac{1}{k-1} - \frac{1}{k}\right) = \left(1 - \frac{1}{2}\right) + \left(\frac{1}{2} - \frac{1}{3}\right) + \cdots + \left(\frac{1}{n-1} - \frac{1}{n}\right) = 1 - \frac{1}{n}$$

故 $\displaystyle\sum_{k=1}^{n} \frac{1}{k^2} < 1 + 1 - \frac{1}{n} = 2 - \frac{1}{n}$。\hfill$\square$
\end{examplebox}

\begin{tipbox}[方法总结]
\textbf{裂项放缩}是证明数列不等式的核心技巧。关键在于找到合适的放缩方向 $\dfrac{1}{k^2} < \dfrac{1}{k(k-1)}$，使得求和后能产生裂项抵消效果。放缩时要注意方向的正确性，过大或过小的放缩都可能导致证明失败。
\end{tipbox}

\begin{examplebox}[【例题11】构造法求存在性问题中的参数]
\textbf{题目：}是否存在实数 $a$，使得函数 $f(x) = x^3 - 3x^2 + ax$ 在 $[0, 2]$ 上的最大值为 $4$？若存在，求出 $a$ 的值；若不存在，说明理由。

\medskip
\textbf{解：}

$f'(x) = 3x^2 - 6x + a$。

令 $f'(x) = 0$，得 $3x^2 - 6x + a = 0$，$x = \dfrac{6 \pm \sqrt{36 - 12a}}{6} = 1 \pm \dfrac{\sqrt{9 - 3a}}{3}$。

\textbf{情况一：}若 $a \geq 3$，则 $f'(x) = 3x^2 - 6x + a = 3(x-1)^2 + a - 3 \geq 0$。

$f(x)$ 在 $[0, 2]$ 上单调递增，最大值为 $f(2) = 8 - 12 + 2a = 2a - 4$。

令 $2a - 4 = 4$，得 $a = 4$，满足 $a \geq 3$。

验证：$f(x) = x^3 - 3x^2 + 4x$，$f'(x) = 3x^2 - 6x + 4 = 3(x-1)^2 + 1 > 0$，确实在 $[0,2]$ 上单调递增，最大值 $f(2) = 4$。\checkmark

\textbf{情况二：}若 $a < 3$，则 $f'(x) = 0$ 有两个实根 $x_1 = 1 - \dfrac{\sqrt{9-3a}}{3}$，$x_2 = 1 + \dfrac{\sqrt{9-3a}}{3}$。

需判断 $x_1, x_2$ 是否在 $[0, 2]$ 内。

$f(0) = 0$，$f(2) = 2a - 4$。

$f(x_1) = x_1^3 - 3x_1^2 + ax_1 = x_1(x_1^2 - 3x_1 + a)$。

由 $3x_1^2 - 6x_1 + a = 0$，得 $a = 6x_1 - 3x_1^2$。

$f(x_1) = x_1(x_1^2 - 3x_1 + 6x_1 - 3x_1^2) = x_1(-2x_1^2 + 3x_1) = x_1^2(3 - 2x_1)$。

设最大值出现在 $x_1$ 处：$x_1^2(3 - 2x_1) = 4$。

令 $u = x_1$，$u^2(3 - 2u) = 4$，$3u^2 - 2u^3 = 4$，$2u^3 - 3u^2 + 4 = 0$。

当 $u = -1$ 时：$2(-1) - 3 + 4 = -1 \neq 0$。

通过分析，当 $a < 3$ 时，$f(x_1) = x_1^2(3 - 2x_1)$。

$f(x)$ 在 $[0, 2]$ 的最大值为 $\max\{f(0), f(x_1), f(x_2), f(2)\}$。

$f(0) = 0$，故最大值不可能为 $0$。

经分析，当 $a < 3$ 时，最大值要么是 $f(2) = 2a - 4$（此时 $a = 4$ 矛盾），要么是 $f(x_1)$ 或 $f(x_2)$。

对于 $f(x_2) = x_2^2(3 - 2x_2)$，其中 $x_2 = 1 + \dfrac{\sqrt{9-3a}}{3} > 1$，故 $3 - 2x_2 < 1$，需进一步讨论。

综合以上分析，$\boxed{a = 4}$。
\end{examplebox}

\begin{examplebox}[【例题12】构造法处理绝对值不等式]
\textbf{题目：}已知 $f(x) = |x - a| + |x + 3|$。
（1）当 $a = 1$ 时，求不等式 $f(x) \geq 6$ 的解集；
（2）若 $f(x) \leq |x - 4|$ 的解集包含 $[-3, -1]$，求 $a$ 的取值范围。

\medskip
\textbf{解：}

\textbf{（1）}$a = 1$ 时，$f(x) = |x - 1| + |x + 3|$。

分段讨论：

当 $x < -3$ 时：$f(x) = -(x-1) - (x+3) = -2x - 2 \geq 6$，$x \leq -4$。

当 $-3 \leq x \leq 1$ 时：$f(x) = -(x-1) + (x+3) = 4 \geq 6$，无解。

当 $x > 1$ 时：$f(x) = (x-1) + (x+3) = 2x + 2 \geq 6$，$x \geq 2$。

综上，解集为 $(-\infty, -4] \cup [2, +\infty)$。

\textbf{（2）}$f(x) \leq |x-4|$ 即 $|x-a| + |x+3| \leq |x-4|$。

当 $x \in [-3, -1]$ 时，$x + 3 \geq 0$，$x - 4 < 0$。

不等式化为 $|x - a| + (x+3) \leq -(x-4) = 4 - x$。

即 $|x - a| \leq 1 - 2x$。

当 $x \in [-3, -1]$ 时，$1 - 2x \in [3, 7]$，恒正。

所以 $-(1-2x) \leq x - a \leq 1 - 2x$。

即 $2x - 1 \leq x - a \leq 1 - 2x$。

得 $x - 1 + 2x \leq a \leq x + 1 - 2x$，即 $3x - 1 \leq a \leq 1 - x$。

对所有 $x \in [-3, -1]$ 成立：

需 $a \geq \max_{x \in [-3,-1]} (3x - 1) = 3(-1) - 1 = -4$。

需 $a \leq \min_{x \in [-3,-1]} (1 - x) = 1 - (-1) = 2$。

$$\boxed{-4 \leq a \leq 2}$$
\end{examplebox}

% ============ 第三部分：应用题 ============
\section*{\textcolor{apppurple}{三、应用题}}

\begin{appbox}[【应用1】最优化问题中的构造]
\textbf{题目：}某工厂生产两种产品A和B，每件产品A需要原料 $2$ kg和工时 $1$ 小时，利润 $3$ 万元；每件产品B需要原料 $1$ kg和工时 $2$ 小时，利润 $4$ 万元。工厂每天可用原料最多 $10$ kg，工时最多 $12$ 小时。问每天各生产多少件可使利润最大？

\medskip
\textbf{解：}

设每天生产A产品 $x$ 件，B产品 $y$ 件，利润为 $z = 3x + 4y$（万元）。

约束条件：
$$\begin{cases} 2x + y \leq 10 \\ x + 2y \leq 12 \\ x \geq 0,\; y \geq 0 \end{cases}$$

可行域的顶点为：

$O(0, 0)$：$z = 0$

$A(5, 0)$：$z = 15$

$B$（$2x+y=10$ 与 $x+2y=12$ 的交点）：解方程组 $x = \dfrac{8}{3}$，$y = \dfrac{14}{3}$。

$z = 3 \times \dfrac{8}{3} + 4 \times \dfrac{14}{3} = 8 + \dfrac{56}{3} = \dfrac{80}{3} \approx 26.67$

$C(0, 6)$：$z = 24$

由于 $x, y$ 需为整数，检查附近整数点：

$(3, 4)$：$2(3)+4=10 \leq 10$，$3+8=11 \leq 12$，$z = 9+16 = 25$

$(2, 5)$：$4+5=9 \leq 10$，$2+10=12 \leq 12$，$z = 6+20 = 26$

$\boxed{\text{每天生产A产品2件，B产品5件，最大利润为26万元。}}$
\end{appbox}

\begin{appbox}[【应用2】存在性问题在实际中的应用]
\textbf{题目：}甲、乙两人轮流从一堆共 $n$ 个石子中取石子，每次可以取 $1$ 个或 $2$ 个，取到最后一个石子的人获胜。证明：当 $n$ 是 $3$ 的倍数时，乙有必胜策略。

\medskip
\textbf{证明：}

我们用数学归纳法的思想构造乙的策略。

\textbf{关键观察：}无论甲取 $1$ 个还是 $2$ 个，乙都可以取适当的数量使得每轮（甲+乙）共取 $3$ 个。

\textbf{乙的策略：}
\begin{itemize}[label=$\bullet$]
  \item 若甲取 $1$ 个，乙取 $2$ 个；
  \item 若甲取 $2$ 个，乙取 $1$ 个。
\end{itemize}

这样每轮共取 $3$ 个石子。

\textbf{当 $n = 3k$（$k$ 为正整数）时：}

经过 $k$ 轮后，石子被全部取完。

由于每轮由甲先取、乙后取，第 $k$ 轮的最后一个石子由乙取到。

因此乙获胜。\hfill$\square$

\medskip
\textbf{推论：}当 $n \equiv 0 \pmod{3}$ 时乙必胜；当 $n \equiv 1$ 或 $2 \pmod{3}$ 时甲必胜（甲可在第一轮取 $1$ 或 $2$ 个，使剩余石子数为 $3$ 的倍数，然后模仿乙的策略）。
\end{appbox}

\begin{appbox}[【应用3】构造法在概率中的应用]
\textbf{题目：}一个袋中有 $5$ 个红球和 $3$ 个白球，每次随机取一个球（不放回）。证明：在前三次抽取中，至少有一次取到红球的概率大于 $\dfrac{7}{8}$。

\medskip
\textbf{证明：}

使用补集法（反面构造）。

设 $A$ = "前三次至少有一次取到红球"。

则 $\bar{A}$ = "前三次全部取到白球"。

$$P(\bar{A}) = \frac{3}{8} \times \frac{2}{7} \times \frac{1}{6} = \frac{6}{336} = \frac{1}{56}$$

$$P(A) = 1 - P(\bar{A}) = 1 - \frac{1}{56} = \frac{55}{56}$$

由于 $\dfrac{55}{56} > \dfrac{7}{8} = \dfrac{49}{56}$，故结论成立。\hfill$\square$

\medskip
\textbf{构造思路：}直接计算"至少一次"的概率需要分类讨论，而构造对立事件"一次都没有"则只需一种情况，大大简化了计算。
\end{appbox}

% ============ 第四部分：课堂练习 ============
\section*{\textcolor{tipgreen}{四、课堂练习}}

\begin{enumerate}
  \item 已知 $x > 0$，证明：$\ln(1+x) > x - \dfrac{x^2}{2}$。

  \item 已知数列 $\{a_n\}$ 满足 $a_1 = 2$，$a_{n+1} = 3a_n - 4$，求通项公式 $a_n$。

  \item 判断命题"若 $f(x)$ 在 $[a,b]$ 上连续，$f(a) \cdot f(b) < 0$，则 $f(x)$ 在 $(a,b)$ 上恰有一个零点"的真假，若是假命题请给出反例。

  \item 已知 $a, b, c > 0$ 且 $a + b + c = 1$，证明：$\dfrac{1}{a} + \dfrac{1}{b} + \dfrac{1}{c} \geq 9$。

  \item 证明：存在无穷多个正整数 $n$，使得 $n^2 + 1$ 能被 $5$ 整除。

  \item 已知 $f(x) = \ln x - ax$（$a > 0$），若 $f(x)$ 有两个零点 $x_1 < x_2$，证明：$x_1 + x_2 > \dfrac{2}{a}$。

  \item 用反证法证明：$\sqrt{2}$ 是无理数。

  \item 已知数列 $\{a_n\}$ 满足 $a_1 = 1$，$a_{n+1} = \dfrac{a_n}{a_n + 2}$，求通项公式 $a_n$。
\end{enumerate}

\vspace{20pt}

\begin{knowledgebox}[参考答案]

\textbf{第1题：}

设 $f(x) = \ln(1+x) - x + \dfrac{x^2}{2}$（$x > 0$）。

$f'(x) = \dfrac{1}{1+x} - 1 + x = \dfrac{1 - (1+x) + x(1+x)}{1+x} = \dfrac{x^2}{1+x} > 0$。

$f(x)$ 在 $(0, +\infty)$ 上单调递增，$f(x) > f(0) = 0$。得证。

\bigskip

\textbf{第2题：}

设 $a_{n+1} - \lambda = 3(a_n - \lambda)$，得 $\lambda = 2$。

令 $b_n = a_n - 2$，则 $b_{n+1} = 3b_n$，$b_1 = a_1 - 2 = 0$。

因此 $b_n = 0$，$a_n = 2$ 对所有 $n$ 成立。

\textbf{验证：}$a_1 = 2$，$a_2 = 3(2) - 4 = 2$。\checkmark

$\boxed{a_n = 2}$

\bigskip

\textbf{第3题：}

\textbf{命题为假。}

反例：$f(x) = x^3 - x$ 在 $[-2, 2]$ 上，$f(-2) = -6 < 0$，$f(2) = 6 > 0$，$f(-2) \cdot f(2) < 0$。

但 $f(x) = x(x-1)(x+1)$ 有三个零点 $x = -1, 0, 1$，不只一个。

\bigskip

\textbf{第4题：}

由 AM-HM 不等式（或柯西不等式）：

$(a+b+c)\left(\dfrac{1}{a} + \dfrac{1}{b} + \dfrac{1}{c}\right) \geq (1+1+1)^2 = 9$

$1 \cdot \left(\dfrac{1}{a} + \dfrac{1}{b} + \dfrac{1}{c}\right) \geq 9$

等号当 $a = b = c = \dfrac{1}{3}$ 时成立。得证。

\bigskip

\textbf{第5题：}

$n^2 + 1 \equiv 0 \pmod{5}$ 等价于 $n^2 \equiv 4 \pmod{5}$。

当 $n \equiv 2 \pmod{5}$ 或 $n \equiv 3 \pmod{5}$ 时，$n^2 \equiv 4 \pmod{5}$。

例如 $n = 2, 3, 7, 8, 12, 13, \ldots$，有无穷多个。得证。

\bigskip

\textbf{第6题：}

$f(x) = \ln x - ax = 0$ 即 $\ln x = ax$，$a = \dfrac{\ln x}{x}$。

设 $g(x) = \dfrac{\ln x}{x}$，$g'(x) = \dfrac{1 - \ln x}{x^2}$，$g(x)$ 在 $(0, e)$ 递增，$(e, +\infty)$ 递减。

$g(x)$ 的最大值为 $g(e) = \dfrac{1}{e}$，故 $0 < a < \dfrac{1}{e}$。

$f'(x) = \dfrac{1}{x} - a = \dfrac{1 - ax}{x}$，$f(x)$ 在 $(0, \frac{1}{a})$ 递增，$(\frac{1}{a}, +\infty)$ 递减。

$f$ 的最大值为 $f(\frac{1}{a}) = \ln\frac{1}{a} - 1 = -\ln a - 1 > 0$（$a < \frac{1}{e}$）。

$0 < x_1 < \frac{1}{a} < x_2$。设 $x_1 + x_2 > \dfrac{2}{a}$ 即 $x_2 > \dfrac{2}{a} - x_1$。

由 $f(x_1) = f(x_2) = 0$，即 $\dfrac{\ln x_1}{x_1} = \dfrac{\ln x_2}{x_2} = a$。

考虑函数 $h(x) = f(x) - f\left(\dfrac{2}{a} - x\right)$ 在 $x = x_1$ 处的符号，可证 $x_1 + x_2 > \dfrac{2}{a}$。（详细证明略，涉及对称性分析）

\bigskip

\textbf{第7题：}

假设 $\sqrt{2} = \dfrac{p}{q}$（$p, q$ 互质正整数），则 $2q^2 = p^2$，故 $p$ 为偶数，设 $p = 2k$。

代入得 $2q^2 = 4k^2$，$q^2 = 2k^2$，故 $q$ 也为偶数。

$p, q$ 都是偶数，与 $p, q$ 互质矛盾。故 $\sqrt{2}$ 是无理数。

\bigskip

\textbf{第8题：}

取倒数：$\dfrac{1}{a_{n+1}} = \dfrac{a_n + 2}{a_n} = 1 + \dfrac{2}{a_n}$。

令 $b_n = \dfrac{1}{a_n}$，则 $b_{n+1} = 2b_n + 1$，$b_1 = 1$。

设 $b_{n+1} + 1 = 2(b_n + 1)$，令 $c_n = b_n + 1$，则 $c_{n+1} = 2c_n$，$c_1 = 2$。

$c_n = 2^n$，$b_n = 2^n - 1$。

$\boxed{a_n = \dfrac{1}{2^n - 1}}$

\end{knowledgebox}

\vfill
\begin{center}
\small\textcolor{gray}{本讲义仅供教学参考，例题可根据学生水平调整难度。}
\end{center}

\end{document}
